\documentclass[conference]{IEEEtran}

% -------------------------------------------------------
% Paquetes
% -------------------------------------------------------
\usepackage[spanish]{babel}
\usepackage[utf8]{inputenc}
\usepackage[T1]{fontenc}
\usepackage{amsmath}
\usepackage{amssymb}
\usepackage{graphicx}
\usepackage{booktabs}
\usepackage{hyperref}
\usepackage{cite}
\usepackage{pgfplots}
\pgfplotsset{compat=1.18}
\usepackage{tikz}
\usepackage{array}
\usepackage{multirow}
\usepackage{url}
\usepackage{fancyhdr}

% -------------------------------------------------------
\begin{document}

\pagestyle{fancy}
\fancyhf{}
\rhead{\includegraphics[height=1.2cm]{uts_logonegro.png}}
\fancyfoot[C]{\thepage}
\renewcommand{\headrulewidth}{0.0pt}
\setlength{\headheight}{1.5cm}

% -------------------------------------------------------
% CABECERA
% -------------------------------------------------------
\title{Boletín | Centro de Inteligencia Artificial e Innovación\\
       \large Mes de octubre de 2026}

\author{
  \IEEEauthorblockN{Miguel Ángel Osorio Orduz}
  \IEEEauthorblockN{Jhoan Sebastian Torres Garcia}
  \IEEEauthorblockN{Frank Camilo Zuñiga Vargas}
  \IEEEauthorblockA{Docente: Carlos Pineda}
  \IEEEauthorblockA{
    Centro de Inteligencia Artificial e Innovación (CIAI)\\
    Colombia\\
    mangelosorio@uts.edu.co\\
    jhoanstorres@uts.edu.co\\
    fczuniga@uts.edu.co
  }
}

\maketitle
\thispagestyle{fancy}

% -------------------------------------------------------
% RESUMEN
% -------------------------------------------------------
\begin{abstract}
Esta primera edición del boletín del CIAI, cubre las novedades
más relevantes en modelos de lenguaje de gran escala durante el mes
de octubre de 2026. Se presentan lanzamientos de OpenAI,
Anthropic, Google, Qwen y DeepSeek, junto con una comparativa de
rendimiento y costo de los modelos más destacados según datos
independientes de Artificial Analysis. También herramientas relevantes, noticias
del mundo de los LLM's.
\end{abstract}

\begin{IEEEkeywords}
LLMs, modelos de lenguaje, herramientas IA, CIAI, benchmarks, novedades
\end{IEEEkeywords}

% ===============================================================
\section{Investigación}
% ===============================================================

Esta sección cubre los modelos y avances publicados durante la semana,
junto con comparativas cuantitativas tomadas de plataformas de evaluación
independientes.

% -------------------------------------------------------
\subsection{Nuevos Modelos}

% -- OpenAI --------------------------------------------------
\subsubsection{OpenAI}

El 26 de octubre, OpenAI lanzó GPT-6.1 Terra en su API junto con
ChatGPT Work y Codex. El modelo reduce los costos a la mitad respecto
a GPT-6.1 Sol, e integra razonamiento visual y uso de computadoras.
Simultáneamente presentaron la \textit{Decisions API} para aplicaciones
responsivas. Previamente, durante el DevDay, habían revelado GPT-6.1
Sol~\cite{gpt61terra}.

El 29 de octubre, OpenAI amplió los límites de uso en los planes
ChatGPT Pro (Pro~100 y Pro~500) y anunció la transición de GPTs
personalizados hacia plugins, la expansión de GPT-6 a todos los planes
y un nuevo enfoque en agentes autónomos mediante la \textit{Agents API}.
También introdujeron nuevas voces para GPT-Live-1.1 y predicciones en
Codex.

En materia de seguridad, el 3 y 6 de octubre se reportaron incidentes
donde agentes autónomos (GPT-6.1 Sol y GPT-6 Astra) evadieron y
desactivaron herramientas de monitoreo interno, negando posteriormente
las acciones realizadas. El 18 de octubre, OpenAI publicó un reporte
detallando la interrupción de más de 20 operaciones cibernéticas globales
e intentos de \textit{spear-phishing} desde China contra sus empleados.
El 28 de octubre se registró una caída global de sus servicios~\cite{openaisecurity}.

% -- Anthropic -----------------------------------------------
\subsubsection{Anthropic}

Anthropic presentó formalmente su formulario S-1 para cotizar en Nasdaq
bajo el símbolo ``ANT'', buscando una valoración de entre \$1.5 y \$2
billones. La compañía reportó ingresos de \$16.286 mil millones en el
primer semestre de 2026, con más de 11.900 clientes corporativos.
Registraron una pérdida contable de \$80.748 mil millones, impulsada
por un cargo único de \$72.4 mil millones por notas convertibles. En el
prospecto, advirtieron formalmente sobre los riesgos catastróficos o
existenciales de la IA~\cite{anthropicipo}.

El 27 y 28 de octubre anunciaron la expansión de su programa
\textit{AI for Science} a 20 nuevas instituciones y revelaron avances
tempranos en el desarrollo de fármacos. También lanzaron Claude
Haiku~5.5, su modelo compacto de mayor rendimiento~\cite{claudehaiku55}.

% -- Google Gemini -------------------------------------------
\subsubsection{Google Gemini}

El 28 de octubre, Google anunció Gemini~4 Carbon, modelo enfocado en
ciberseguridad y tareas prolongadas de programación, con acceso
restringido a defensores cibernéticos y socios del programa Fairwind.
El 13 de octubre se habilitó el acceso general a Gemini~4 Argon para
suscriptores y desarrolladores, integrando el sistema \textit{Antigravity}
en la app de macOS~\cite{gemini4carbon}.

El 30 de octubre, Google introdujo la edición de documentos (Docs,
Sheets, Code, PDF) a pantalla completa directamente en el chat. La
función \textit{Gems} comenzó su transición para convertirse en
\textit{Skills} en noviembre, sumando integración con más de 40
aplicaciones~\cite{geminieco}.

% -- Qwen ----------------------------------------------------
\subsubsection{Qwen}

El 27 de octubre, Alibaba lanzó \texttt{qwen-audio-3.1-realtime-flash},
modelo de voz bidireccional con 262K tokens de contexto, búsqueda web
y clonación de voz. El 2 y 9 de octubre publicaron las versiones
\textit{pro} y \textit{turbo} de su modelo de generación visual
Qwen-Image-2.1~\cite{qwenaudio}.

Adicionalmente, se reportaron pruebas en la plataforma Arena del modelo
``Kalmia'', que se especula es la versión preliminar de Qwen~4.

% -- DeepSeek ------------------------------------------------
\subsubsection{DeepSeek}

DeepSeek levantó \$15 mil millones con apoyo de Tencent y CATL,
alcanzando una valoración pre-money de \$74 mil millones, y aceleró
sus planes para una IPO en el mercado STAR de Shanghái para 2027.
Esta financiación ocurrió días después de que el gobierno de EE.~UU.
ordenara a agencias federales civiles eliminar la tecnología de DeepSeek
de sus infraestructuras~\cite{deepseekfunding}.

El 17 de octubre actualizaron DeepSeek V3 con mayor velocidad y menor
cómputo. Posteriormente lanzaron DeepSeek-V4.1-Flash, que integra
comprensión visual nativa y reemplaza a la generación V4. Su API ahora
soporta el formato nativo de Responses de OpenAI~\cite{deepseekv4}.

% -------------------------------------------------------
\subsection{Analítica de Rendimiento}

Datos de costo y latencia tomados de Artificial~Analysis para los modelos
más relevantes lanzados o actualizados esta semana~\cite{artificialanalysis}.
El TTFT corresponde al modo de esfuerzo estándar o bajo según disponibilidad.
Gemini~4 Carbon no tiene precio público de API (acceso restringido Fairwind).

\begin{table}[h]
\centering
\caption{Comparativa de Costo y Latencia -- Octubre 2026}
\label{tab:comparativa}
\begin{tabular}{@{}lcccc@{}}
\toprule
\textbf{Modelo} & \textbf{Input} & \textbf{Output} & \textbf{TTFT} & \textbf{Tok/s} \\
                & (\$/1M tok.)   & (\$/1M tok.)    & (s)           &                \\
\midrule
GPT-6.1 Terra      & 1.00  & 5.00  & 1.29  & 91  \\
Claude Haiku 5.5   & 0.10  & 0.50  & 1.22  & 240 \\
Gemini 4 Argon     & 2.00  & 10.00 & N/D   & N/D \\
Gemini 4 Carbon    & --*   & --*   & 97.48 & 99.8 \\
DeepSeek V4.1 Flash & 0.15 & 0.60  & N/D   & N/D \\
Qwen3.8 Max        & 2.00  & 6.00  & 1.95  & 54  \\
\bottomrule
\multicolumn{5}{@{}l}{\footnotesize *Acceso restringido; sin precio público.} \\
\multicolumn{5}{@{}l}{\footnotesize Gemini 4 Argon: precio introductorio hasta nov.~30, 2026.} \\
\multicolumn{5}{@{}l}{\footnotesize DeepSeek V4.1 Flash: tarifa off-peak; peak \$0.30/\$1.20.} \\
\end{tabular}
\end{table}

\begin{figure}[h]
\centering
\begin{tikzpicture}
\begin{axis}[
  width=\columnwidth,
  height=5.5cm,
  ybar,
  symbolic x coords={
    GPT-6.1 Terra,
    Haiku 5.5,
    Gemini 4 Argon,
    DS V4.1 Flash,
    Qwen3.8 Max
  },
  xtick=data,
  ylabel={Costo input (\$/1M tok.)},
  ymin=0,
  bar width=0.35cm,
  nodes near coords,
  nodes near coords align={vertical},
  enlarge x limits=0.15,
  x tick label style={font=\tiny, rotate=30, anchor=east},
]
\addplot coordinates {
  (GPT-6.1 Terra,    1.00)
  (Haiku 5.5,        0.10)
  (Gemini 4 Argon,   2.00)
  (DS V4.1 Flash,    0.15)
  (Qwen3.8 Max,      2.00)
};
\end{axis}
\end{tikzpicture}
\caption{Costo de entrada por millón de tokens (modelos con precio
público disponible). Fuente: Artificial~Analysis, oct.~2026.}
\label{fig:costo}
\end{figure}

\begin{figure}[h]
\centering
\begin{tikzpicture}
\begin{axis}[
  width=\columnwidth,
  height=5.5cm,
  ybar,
  symbolic x coords={
    GPT-6.1 Terra,
    Haiku 5.5,
    Gemini 4 Carbon,
    Qwen3.8 Max
  },
  xtick=data,
  ylabel={TTFT (s)},
  ymin=0,
  bar width=0.35cm,
  nodes near coords,
  nodes near coords align={vertical},
  enlarge x limits=0.15,
  x tick label style={font=\tiny, rotate=30, anchor=east},
]
\addplot coordinates {
  (GPT-6.1 Terra,    1.29)
  (Haiku 5.5,        1.22)
  (Gemini 4 Carbon, 97.48)
  (Qwen3.8 Max,      1.95)
};
\end{axis}
\end{tikzpicture}
\caption{Time to First Token (TTFT) en segundos, esfuerzo estándar/bajo.
Gemini~4 Carbon opera en modo razonamiento extendido, lo que explica su
TTFT elevado. Fuente: Artificial~Analysis, oct.~2026.}
\label{fig:ttft}
\end{figure}

% ===============================================================
\section{Herramientas}
% ===============================================================

Herramientas, extensiones y entornos de creación publicados o
actualizados este mes. Enfoque en opciones gratuitas y no comerciales.

% -------------------------------------------------------
\subsection{OpenCut}

OpenCut es un editor de video de código abierto bajo licencia MIT,
diseñado como alternativa directa a CapCut. Corre enteramente en el
navegador mediante WebAssembly, sin requerir cuenta ni subir archivos
a servidores externos. Su arquitectura está basada en Next.js para la
web, con un núcleo compartido en Rust para composición GPU, efectos y
máscaras. En mayo de 2026 el equipo inició una reescritura completa
con arquitectura orientada a plugins y una API headless planificada.
A octubre de 2026 el proyecto acumula más de 93.000 estrellas en GitHub
y se encuentra en estado pre-lanzamiento activo~\cite{opencut}.

\begin{itemize}
  \item \textbf{Tipo:} Aplicación web / desktop (en desarrollo)
  \item \textbf{Licencia:} MIT
  \item \textbf{Web:} \url{https://opencut.app}
  \item \textbf{Repositorio:} \url{https://github.com/OpenCut-app/OpenCut}
\end{itemize}

% -------------------------------------------------------
\subsection{Lemon Studio}

Lemon Studio es una estación de trabajo de audio digital (DAW)
de código abierto, escrita completamente en Rust. Está diseñada
como alternativa libre a FL Studio, con la que mantiene
interoperabilidad directa: abre y guarda proyectos \texttt{.flp}
de FL Studio 24.2 en adelante de forma no destructiva. Incluye
secuenciador por pasos, piano roll completo, mezclador con ruteo,
y soporte para plugins CLAP y VST3. Corre en Windows, Linux y
macOS. El proyecto se encuentra en estado pre-alpha; aún no
distribuye binarios compilados y debe construirse desde el código
fuente~\cite{lemonstudio}.

\begin{itemize}
  \item \textbf{Tipo:} Aplicación desktop (DAW)
  \item \textbf{Licencia:} GPL-3.0
  \item \textbf{Repositorio:} \url{https://github.com/ItsAlphaNeon/Lemon-Studio}
\end{itemize}

% -------------------------------------------------------
\subsection{ArtCraft}

ArtCraft es un ecosistema de herramientas de escritorio de código
abierto construidas en Rust, desarrollado por el equipo de Storytold.
Se divide en dos productos principales~\cite{artcraft}:

\textbf{Crafting Apps:} trece aplicaciones nativas para flujos de
trabajo creativos y de oficina, instalables mediante el ArtCraft
Launcher. Cubren edición de imágenes (PhotoCraft), ilustración
vectorial (VectorCraft), edición de video (FilmCraft), ajuste de
color (LightCraft), PDF (PdfCraft), efectos visuales (EffectCraft),
diseño gráfico (DesignCraft), diseño de interfaces (FigCraft),
procesamiento de texto (WordCraft), audio (SoundCraft), presentaciones
(DeckCraft), modelado CAD (CADCraft) y hojas de cálculo (GridCraft).
Son equivalentes funcionales de las suites de Adobe y Microsoft
Office, sin suscripción ni costo.

\textbf{ArtCraft Studio:} aplicación de generación de imágenes y
video con IA. Permite componer escenas en un canvas 2D/3D antes de
invocar el modelo generativo, dando control real sobre posición de
personajes, perspectiva, capas y assets mixtos. Soporta modelos como
Seedance 2.5, Nano Banana 2, Kling y otros. Disponible como app de
escritorio y en el navegador sin instalación.

Ambos productos son gratuitos, sin suscripción y de código abierto.

\begin{itemize}
  \item \textbf{Tipo:} Suite de aplicaciones desktop + estudio IA generativa
  \item \textbf{Licencia:} MIT / Apache 2.0
  \item \textbf{Web:} \url{https://getartcraft.com}
  \item \textbf{Launcher:} \url{https://github.com/storytold/craft-launcher}
  \item \textbf{Studio:} \url{https://github.com/storytold/artcraft}
  \item \textbf{App web:} \url{https://app.getartcraft.com}
\end{itemize}

% -------------------------------------------------------
\subsection{OpenIDEA}

OpenIDEA es un entorno de desarrollo integrado (IDE) de código abierto
construido en Rust y TypeScript sobre Tauri, diseñado como alternativa
nativa y ligera a IntelliJ IDEA. Adopta los atajos de teclado, ventanas
de herramientas y la filosofía de búsqueda global de los IDEs de
JetBrains, sin afiliación con esa empresa~\cite{openidea}.

Integra editor, inteligencia de código (LSP), depurador, herramientas
de Git, terminal, cliente de bases de datos y cliente HTTP nativo.
Soporta más de 60 lenguajes de programación y cuenta con versiones para
macOS, Windows y Linux.

En rendimiento, abriendo un proyecto Python de 770 archivos frente a
IntelliJ IDEA 2026.2.3: carga la ventana en 0.9s (vs. 3.1s), consume
169~MB de RAM (vs. más de 2000~MB) y estabiliza el CPU en 1.1s
(vs. más de 13s). La versión más reciente es la v0.1.28.

\begin{itemize}
  \item \textbf{Tipo:} IDE de escritorio
  \item \textbf{Licencia:} Apache 2.0
  \item \textbf{Repositorio:} \url{https://github.com/l0ng-ai/openidea}
\end{itemize}

% ===============================================================
\section{Proyectos CIAI}
% ===============================================================

Avances, publicaciones y actividades del Centro de Inteligencia
Artificial e Innovación durante la semana.

% -------------------------------------------------------
\subsection{[Nombre del Proyecto 1]}

[Descripción del avance, estado actual, próximos pasos y personas
involucradas.]

% ===============================================================
\section{Noticias}
% ===============================================================

Eventos relevantes de la industria de la
Ingeniería de sistemas y afines.

% -------------------------------------------------------
\subsection{REA: Reverse Engineer Anything alcanza 234.000 estrellas en GitHub}

REA (\textit{Reverse Engineer Anything}) es un proyecto de código
abierto bajo licencia MIT que permite a un agente de IA, como Claude
Code, Codex o Cursor, analizar binarios compilados sin acceso al
código fuente. El repositorio \texttt{morluto/rea} fue creado el 14
de abril de 2026 y en octubre superó las 234.000 estrellas y 58.000
forks, convirtiéndose en uno de los proyectos de mayor crecimiento
en la historia reciente de GitHub~\cite{rea}.
Este proyecto se volvió ultra famoso en octubre de 2026, como
ayuda para crear aplicaciones con Ingeniería Inversa.

El proyecto no es un decompilador: utiliza al agente de IA para
reconstruir el comportamiento de funciones a partir de instrucciones,
cadenas de texto y patrones observables en el binario. Se instala
vía \texttt{npx rea-agents setup} y requiere Node.js 22.19 o superior.

Durante la semana del 13 de octubre se detectó un repositorio
falso bajo la organización \texttt{rea-tools} que copiaba el README
y distribuía un troyano en un ejecutable \texttt{.exe} para Windows.
El README oficial fue actualizado para advertir que REA debe obtenerse
únicamente desde \texttt{github.com/morluto/rea}, el sitio
\texttt{rea.tools} o el paquete npm \texttt{rea-agents}~\cite{rea}.

\begin{itemize}
  \item \textbf{Repositorio oficial:} \url{https://github.com/morluto/rea}
  \item \textbf{Sitio:} \url{https://rea.tools}
  \item \textbf{Paquete npm:} \texttt{rea-agents}
\end{itemize}

% -------------------------------------------------------
\subsection{OpenIDEA entra en GitHub Trending con rendimiento
superior a IntelliJ}

OpenIDEA acumuló más de 3.500 estrellas y 290 forks en GitHub en
pocas semanas, impulsado por benchmarks que muestran una reducción
del 91\,\% en consumo de RAM y tiempos de carga hasta 3.4 veces
más rápidos frente a IntelliJ IDEA 2026.2.3~\cite{openidea}. El
proyecto se encuentra en desarrollo activo: los commits más recientes
incluyen soporte para arrastrar pestañas como ventanas independientes,
espacios de trabajo multiproyecto, mejoras en LSP para auto-importaciones
de archivos no abiertos, y soporte de sintaxis mejorado para Angular,
Kotlin, Scala y Dart.

% -------------------------------------------------------
\subsection{Character.AI se acoge al Capítulo 11 de bancarrota}

El 12 de octubre de 2026, Character Technologies, empresa detrás de
Character.AI, se acogió voluntariamente al Capítulo~11 de la ley de
bancarrotas en Delaware. Su director ejecutivo atribuyó la situación
a ``importantes vientos en contra legales, regulatorios y operativos''.
La compañía lleva a cabo un proceso de venta supervisado por el
tribunal, con cierre esperado entre finales de noviembre y principios
de diciembre de 2026~\cite{characterai}.

A pesar del procedimiento legal, la plataforma permanece en línea.
Los usuarios conservan la opción de eliminar sus cuentas o exportar
su información desde la app o el sitio web. Entre sus funciones
recientes destaca \textit{Breaking News Character}, un presentador
impulsado por IA para titulares y comentarios sobre eventos actuales.

% -------------------------------------------------------
\bibliographystyle{IEEEtran}
\bibliography{referencias}

\end{document}
